\documentclass[journal]{IEEEtran}

% Course template packages, plus the packages used by the existing report.
\usepackage{graphicx}
\usepackage[margin=1in]{geometry}
\usepackage{amsmath,amsthm,amssymb}
\usepackage{amsfonts}
\usepackage{float}
\usepackage{booktabs}
\usepackage[T1]{fontenc}
\usepackage[utf8]{inputenc}
\usepackage{enumitem}
\usepackage[hidelinks]{hyperref}
\usepackage{xcolor}

\newcommand{\system}{\textsc{ProvGate}}
\newcommand{\attack}{\textsc{Attack}}
\newcommand{\placebo}{\textsc{Placebo}}
\newcommand{\clean}{\textsc{Clean}}
\newcommand{\fullarm}{\textsc{Full}}
\newcommand{\caparm}{\textsc{Capability-Only}}
\newcommand{\promptcaparm}{\textsc{Prompt+Capability}}
\newcommand{\allowarm}{\textsc{Allow-All}}

\hyphenation{op-tical net-works semi-conduc-tor Prove-nance}

\begin{document}
\pagenumbering{arabic}
\def\arraystretch{1.2}

\title{Provenance-Aware Capability Gating against\\
Indirect Prompt Injection in Multi-Agent Systems}
\author{Chi Kit WONG, Haotian WANG, Guanrun PENG, Boyong HOU\\
\{ckwong627, hwang418, gpeng615, bhou204\}@connect.hkust-gz.edu.cn}
\date{Summer 2026}

\maketitle

\begin{abstract}
Tool-using agents can confuse data supplied by an attacker with authority
granted by the user.  We study the smallest auditable version of this problem:
ordinary capability gating checks whether a value is on an allow-list, whereas
PACT (Provenance-Aware Capability Tracking/Control) also checks whether that
value is trusted for its parameter role.  We implement a deterministic PACT
reference monitor and evaluate four frozen cases.  PACT allows a user-selected
Alice recipient, denies an externally supplied Bob even when Bob is on the same
allow-list, allows external text at a low-risk content role, and allows a user
value after an explicitly registered Base64 transformation.  Every case records
the argument role, value provenance, both decisions, execution and side-effect
count, rejection reason, and SHA-256 decision/provenance evidence.  The public
AgentDojo benchmark is retained only as an external attack-realism baseline; it
does not directly evaluate PACT.  The result is a small, reproducible mechanism
demonstration rather than a claim of semantic information-flow security.
\end{abstract}

\IEEEpeerreviewmaketitle

\section{Introduction}

Large language models increasingly act through tools rather than only producing
text.  Retrieval makes such systems useful, but it also introduces a confused-
deputy problem: an attacker can place an instruction in an email or document,
and a model may execute it using authority granted by the user.  This is an
\emph{indirect prompt injection} because the adversarial instruction arrives
through data selected at run time rather than through the user message
\cite{greshake2023indirect}.  Evaluations of tool-using agents show that neither
ordinary helpfulness prompting nor simple refusal instructions reliably
separate data from commands \cite{debenedetti2024agentdojo,zverev2024separate}.

This project asks:

\begin{quote}
Can a provenance-aware role check prevent an externally supplied allow-listed
value from laundering authority into a high-impact parameter, while still
allowing benign user values, low-risk external content, and explicitly
registered transformations?
\end{quote}

We deliberately make a narrow claim.  The prototype is a local synthetic office,
the capability oracle is assumed correct, and leakage tracking recognizes exact
registered values.  It is an auditable experiment in enforcement mechanisms,
not a proof of general language-model security.

\section{Background \& Motivation}

Prompt-level instruction hierarchy attempts to teach a model that trusted system
and user instructions outrank untrusted content \cite{wallace2024hierarchy}.
This can reduce attacks but leaves the same probabilistic component responsible
for interpreting both data and policy.  Capability systems instead provide an
unforgeable, least-authority description of the operations a component may
perform \cite{dennis1966capabilities,saltzer1975protection}.  Information-flow
control tracks where information came from and constrains where it may go
\cite{denning1976lattice}.

\system{} combines these ideas at the tool boundary.  The model may still read
an attack and may still propose a malicious call.  Security rests on a small,
deterministic reference monitor that checks the call against task authority and
against the provenance and sensitivity of outbound values.  The distinction
between proposal, policy decision, execution, and world mutation is preserved in
an append-only audit trail.

\section{Target System Description}

\subsection{Agents and mock tools}

The Reader Agent may call \texttt{search\_emails} and
\texttt{read\_email}.  It returns a task summary plus exact evidence and its
event identifiers.  The Action Agent receives the original request, the summary,
and retrieved evidence.  Depending on the scenario, it may propose
\texttt{read\_file}, \texttt{search\_calendar},
\texttt{send\_email}, or \texttt{create\_calendar\_event}.  All resources and
side effects live in a fresh in-memory world for each attempt.

The separation is intentional: it creates a realistic propagation path from an
untrusted retrieval agent to a privileged action agent while retaining complete
evidence about exposure.  A structured context event links the exact resource-
read event to the Action Agent's context.

\subsection{Capabilities}

Each task supplies a finite capability set.  A capability binds a tool to zero
or more resource identifiers, recipient allow-lists, parameter bounds, maximum
call counts, and maximum outbound sensitivity.  For example, a reply task may
authorize one message to \texttt{alice@example.com} while excluding added carbon
copy recipients.  A scheduling task may constrain participants, start/end
windows, and duration.  Calls outside this envelope are denied even if a model
describes them as mandatory.

\subsection{Provenance and sensitivity}

Every relevant value carries an authority label (trusted, untrusted, or
derived), a sensitivity label (public, internal, or secret), and explicit source
and parent event identifiers.  Derived values conservatively retain their
sources.  The prototype additionally registers exact synthetic protected
literals.  If such a literal flows from a protected resource into an outbound
argument, the gateway can deny the call even when the tool and recipient are
otherwise authorized.  The four-case PACT experiment below isolates this
source-and-role check from ordinary allow-list checking.  Earlier T5--T6 and
complete-factorial artifacts are retained as historical supplementary results;
they are not part of the primary PACT table.

\section{Threat Model}

\paragraph{Adversary.}
For the broader prototype, the attacker controls one external-content block in
a synthetic email consumed by the Reader Agent.  It can issue natural-language
instructions designed to cause an unauthorized recipient, file read, calendar
mutation, or disclosure.  The primary PACT artifact does not run a live model:
it uses four deterministic synthetic calls with fixed provenance labels.  The
formal Red Agent and frozen candidate payloads belong only to the retained
supplementary attack-corpus experiments.

\paragraph{Trusted computing base.}
The scenario manifest, capability gateway, provenance tracker, mock tools,
append-only artifact writer, and independent evaluator are trusted.  The Reader
and Action model outputs, retrieved email content, summaries derived from it, and
all tool arguments proposed by a model are untrusted.  We assume the manifest
correctly captures the user's intended authority.

\paragraph{Attacker goals.}
Success requires an observable world effect: a forbidden message recipient, an
unauthorized calendar event or resource access, or an exact registered synthetic
secret in an outbound message.  Model text that merely repeats an instruction is
not counted as an executed effect.  A denied proposal is an attempted and
blocked attack, not a successful attack.

\paragraph{Out of scope.}
The system exposes no shell, payment, deletion, real account, credential, or
external network tool.  We do not claim protection against a compromised
gateway, an incorrect capability manifest, semantic paraphrases of a secret,
side channels, encoded values, or attacks that exploit the model-serving stack.

\section{Attack \& Defense Design}

\subsection{Attacks}

The frozen main experiment contains four deterministic calls to the same
synthetic \texttt{send\_email} tool.  The calls model the two relevant ways in
which an indirect prompt injection can influence a tool invocation: it can
replace a high-impact argument (the recipient/target) with an externally
supplied value, or it can supply ordinary text that should remain usable as
message content.  The fourth case is a control for a legitimate user value
that is changed by an explicitly registered format transformation.

Each argument carries a role and an immutable provenance label.  The central
attack is authority laundering: Bob is present in the capability allow-list,
but the value is introduced by an external email even though the user asked
to contact Alice.  This separates a value-level allow-list check from a
source-and-role check without relying on a model's refusal wording.

\begin{table}[t]
\centering
\caption{Frozen PACT mechanism cases.}
\label{tab:pact-cases}
\resizebox{\columnwidth}{!}{%
\begin{tabular}{llll}
\toprule
Case & Capability-only & PACT & Side effects \\
\midrule
User Alice $\rightarrow$ recipient & Allow & Allow & 1 \\
External Bob, Bob allow-listed & Allow & Deny & 0 \\
External text $\rightarrow$ content & Allow & Allow & 1 \\
Registered Base64 user transform & Allow & Allow & 1 \\
\bottomrule
\end{tabular}}
\end{table}

\subsection{Defenses}

Ordinary Capability-only checks compare a proposed value against the
capability allow-list but ignore its provenance.  PACT first performs that
same value check and then requires a trusted provenance level for the
high-impact roles \texttt{recipient}, \texttt{target}, and \texttt{control}.
External values remain permitted for the low-risk \texttt{content} role.  A
format change is trusted only when the transformation registry contains the
exact source hash, output hash, transformation name, and trusted source
authority; the policy therefore uses an auditable provenance chain rather
than fuzzy string matching.

The second row is the authority-laundering/confused-deputy case: Bob is a legal
value according to the manifest, but an external email is not trusted to bind
Bob to the recipient role.  The third row prevents an over-broad policy from
blocking all retrieved data.  The fourth row is accepted only when the exact
source hash, output hash, transform name, and trusted source authority are
present in an append-only transformation registry.

\section{Evaluation}

\subsection{Experimental Setup}

The experiment has four frozen cases and two policy conditions.  Every case is
executed with a fresh deterministic gateway and a synthetic side-effect ledger.
The output records the parameter name, role, value, provenance source and
transform chain, Capability-only decision, PACT decision, execution flag,
side-effect count, reason, provenance digest, and decision-log SHA-256.

\begin{table}[t]
\centering
\scriptsize
\caption{Minimum deterministic PACT evaluation.}
\label{tab:minimum-matrix}
\begin{tabular}{p{0.33\columnwidth}ccr}
\toprule
Case & Capability-only & PACT & Side effects \\
\midrule
User selects Alice & Allow & Allow & 1 \\
External Bob (allow-listed) & Allow & Deny & 0 \\
External content body & Allow & Allow & 1 \\
Registered Base64 transform & Allow & Allow & 1 \\
\bottomrule
\end{tabular}
\end{table}

The reproducible entry point is
\texttt{scripts/run\_pact\_minimum.py}.  It refuses to overwrite an existing
artifact directory.  The committed artifact contains \texttt{results.csv},
\texttt{results.json}, \texttt{decision\_log.jsonl}, and
\texttt{architecture.mmd}, together with a manifest binding all output
SHA-256 values to the code commit.

The enforcement path is: a user or external email supplies a value; the
gateway attaches its provenance; the value is bound to a parameter role; and
the capability and PACT checks jointly decide whether the tool call can reach
the side-effect ledger.  In the authority-laundering case, the capability
check returns Allow while the PACT role check returns Deny, so no tool call and
no side effect are recorded.

\begin{figure}[t]
\centering
\fbox{%
\begin{minipage}{0.86\columnwidth}
\centering
\texttt{user / external email}\par
$\downarrow$ attach provenance\par
\texttt{recipient | target | content}\par
$\downarrow$ capability + PACT checks\par
\texttt{ALLOW: tool call + side effect}\par
\texttt{DENY: no call + no side effect}
\end{minipage}}
\caption{The deterministic PACT enforcement path at the tool boundary.}
\label{fig:pact-flow}
\end{figure}

\paragraph{Decision-log excerpt.}
A representative blocked row records
\texttt{case\_id=external-bob-allowlisted}, \texttt{role=recipient}, and
\texttt{source=external:email-evil}.  Its decisions are
\texttt{capability=Allow}, \texttt{pact=Deny}, \texttt{executed=0}, and
\texttt{side\_effects=0}; the complete row also stores the provenance
digest and decision-log SHA-256 in the committed JSONL artifact.

\subsection{Outcome Definitions}

For each argument, $C$ denotes the ordinary Capability-only decision and $P$
denotes the PACT decision.  $X$ denotes an executed tool call and $S$ denotes
a synthetic side effect.  The mechanism checks require $X=S=1$ only when
$C=P=1$.  A denied call must have $X=S=0$.  A provenance digest and a decision
log hash are required for every row; these hashes provide evidence binding the
recorded value and source label to the decision.

\subsection{Results}

The four cases reproduce exactly as follows:

\begin{table}[t]
\centering
\caption{PACT versus ordinary capability gating.}
\label{tab:pact-results}
\begin{tabular}{p{0.38\columnwidth}ccr}
\toprule
Case & $C$ & $P$ & $S$ \\
\midrule
User Alice as recipient & Allow & Allow & 1 \\
External Bob as recipient & Allow & Deny & 0 \\
External email as content & Allow & Allow & 1 \\
Registered Base64 user value & Allow & Allow & 1 \\
\bottomrule
\end{tabular}
\end{table}

Capability-only allows all four capability-valid values (4/4).  PACT allows
the user-selected value, the low-risk external content, and the registered
transformation (3/4), while blocking the externally supplied allow-listed Bob
in the high-trust recipient role (1/1 targeted authority-laundering case).
The blocked row has zero tool calls and zero synthetic side effects; the three
allowed rows each have one side effect.

The second case is the central result.  Capability-only cannot distinguish an
allow-listed Bob supplied by the user from an allow-listed Bob supplied by an
external email.  PACT rejects only the high-trust role binding, without
blocking the content case.  The registered transformation case demonstrates
that provenance can survive a format change when the transformation is
explicitly recorded and hash-verified.

\subsection{External AgentDojo Baseline}

We retain the pinned Qwen3-VL-8B AgentDojo result as an external realism check:
Qwen3-VL-8B remains vulnerable to indirect prompt injection, and the effect of
a simple repeat-user-prompt defense depends on attack family and suite.  This
result motivates the problem but does not evaluate PACT and is not merged with
Table~\ref{tab:pact-results}.  The frozen baseline and its SHA-256 manifests
remain in the separate external-baseline artifact directories; their exact
paths are listed in the appendix.

\section{Analysis and Discussion}

\subsection{Interpretation}

The results support a precise mechanism claim rather than a broad security
claim.  Ordinary capability gating answers ``is this value allowed?''  PACT
adds ``is this source trusted for this role?''  Thus the allow-listed external
Bob is denied only when bound to \texttt{recipient}; external email text still
reaches \texttt{content}, and a user value survives a registered Base64
transformation.  The four cases are deterministic and intentionally small so
that every decision and side effect can be inspected.

\subsection{Limitations}

The capability manifest is assumed correct and the PACT cases use a synthetic
side-effect ledger rather than a production messaging service.  Provenance is
explicit and hash-based, not semantic: unregistered decoding, paraphrases,
partial values, and transformations omitted from the registry are outside the
claim.  The experiment uses one tool contract, one local implementation, one
model-independent deterministic policy, and no multi-seed or multi-model
variance.  The AgentDojo benchmark is retained as an external baseline only.
Secret Broker, complete L3 user confirmation, semantic provenance, adaptive
attackers, and broader application suites are future work.

\subsection{Ethics and Safety}

All names, email addresses, files, calendar entries, secrets, and side effects
are synthetic and local.  The model receives no shell, real messaging, real
calendar, payment, deletion, credential, or arbitrary network capability.
Payload generation is separated from formal execution, artifacts are hashed,
and the public conclusions state the defense assumptions and failure modes.  The
work aims to improve defensive measurement.  Attack examples should be shared
with enough context to reproduce the benchmark but not represented as a
guaranteed bypass of unrelated production systems.

\subsection{Reproducibility}

The primary artifact is self-contained.  The directory
\texttt{pact-minimum-v2} contains the four-case \texttt{results.json} and
\texttt{results.csv}, the append-only \texttt{decision\_log.jsonl}, the
Mermaid architecture file, and a manifest binding output SHA-256 values to the
recorded code commit.  The deterministic replay entry point is
\texttt{scripts/run\_pact\_minimum.py}; it refuses to overwrite an existing
artifact directory.  Unit tests run in the
cluster's \texttt{test} Conda environment and require no model server.  The
pinned AgentDojo artifacts are separate external-baseline evidence.

\subsection{AI-Use Disclosure}

Generative AI systems were used as engineering and writing assistants: to help
scaffold code and tests, review threat-model and experimental-design choices,
draft documentation, and, for historical supplementary experiments, generate
offline synthetic attack candidates.  Those historical corpus records include
the model, prompt, seed, raw output, normalized output, and hash.  The primary
PACT four-case artifact is deterministic and does not depend on a model
response.  Human team members remain responsible for the research question,
security assumptions, scenario authorization, code review, executed
experiments, statistical interpretation, source verification, and final prose.
AI-generated content is not accepted as evidence; claims are supported by saved
artifacts or cited literature.

\section{Peer Evaluation}

\subsection{Individual Contribution Statements}

\begin{itemize}[nosep]
  \item Chi Kit WONG: Led project scoping and security-experiment design; developed the threat model and provenance-aware capability-gating architecture; designed and froze the original and hardened offline attack corpora; specified the $2 \times 2$ ablation and deterministic sink-pressure tests; coordinated experiment execution and artifact validation; maintained the GitHub/Overleaf reproducibility package and integrated the final report.
  \item Haotian WANG:
  \item Guanrun PENG:
  \item Boyong HOU:
\end{itemize}

\paragraph{Equal contribution.}
All four members made equal contributions.  The statements above describe
individual roles and do not imply unequal overall contribution; the remaining
statements will be completed after final team agreement.

\subsection{Teammate Evaluation Scores \& Justifications}

The team reports equal contribution for all four members.  No differentiated
scores or justifications are asserted in this report; any separate course
peer-evaluation form should use the same equal-contribution statement.

\section{Conclusion}

We presented PACT, a deterministic provenance-aware capability monitor for
multi-agent tool calls.  In four frozen cases, ordinary capability checks allow
both user- and externally supplied allow-listed values, whereas PACT rejects
the externally supplied value when it is bound to the high-trust recipient
role.  PACT still allows external content at a low-risk role and allows a user
value after a registered, hash-verified Base64 transformation.  The artifact
contains the decision records, side-effect counts, provenance digests, and
manifest hashes needed to reproduce these claims.  The result is a focused
course-project mechanism demonstration, not a claim of complete semantic
information-flow security.

\nocite{*}
\bibliographystyle{IEEEtran}
\bibliography{references}

\newpage
\section*{Appendices}

\subsection*{Artifact and replay pointers}

The primary PACT artifact is under \path{artifacts/pact-minimum-v2/}; its
manifest records output hashes and the code commit.  Historical original and
hardened factorial plans, frozen attack corpora, and AgentDojo summaries remain
under their existing paths as supplementary evidence and are not mixed with
the four-case PACT results.  The report can be rebuilt with
\path{bash scripts/compile_report.sh}; deterministic controller tests use the
repository's \texttt{test} Conda environment.

\end{document}
